\subsection{Source and Coverage}

We use the \emph{Electric Vehicle Charging Station Data} dataset
published by the City of Boulder, Colorado on its Open Data
portal~\cite{boulderdata}. The dataset is updated regularly; the
snapshot we work with covers the period from August 22, 2018 to
November 30, 2023 (1\,927 days), and contains 148\,136 individual
charging sessions across 26 city-owned level-2 stations. Each session
record carries 15 fields, of which the most relevant for forecasting
are the station identifier, the start and end timestamps, the
charging duration, and the energy delivered in kWh. There are
essentially no missing values: only four sessions out of 148\,136 have
a null end-time.

\subsection{Station and Aggregation}

We focus on the station with the highest session volume,
N~Boulder Rec~1, with 16\,278 sessions in the period of interest. The
reason for picking a single station is simple: with one station the
daily series is dense and continuous, with only 5.2\% of days at
zero. Aggregating across all 26 stations would smooth out the
station-level structure we want to learn (commute patterns at
recreational vs.\ commercial locations are very different), and
working with multiple separate series would require multi-task models
that go beyond the scope of this project. The sessions for
N~Boulder Rec~1 are summed by date to produce a univariate daily
series of 1\,927 observations, with mean 74.6~kWh/day, standard
deviation 50.3~kWh, and a range of 0--310~kWh. 100 days have zero
energy (mostly holidays or short station outages); we preserve these
zeros instead of imputing them, since they reflect real operational
patterns that the model should learn.

\begin{figure}[!htbp]
\centering
\includegraphics[width=0.95\textwidth]{fig_series}
\caption{Daily energy consumption at N~Boulder Rec~1 from August 2018
to November 2023. The orange band marks the 30-day held-out test
window; everything to the left of it is training data. The COVID-19
drop in 2020 and the post-pandemic recovery are clearly visible.}
\label{fig:series}
\end{figure}

\subsection{Exogenous Variables}

For the models that support them (SARIMAX, ARIMAX, LSTM, PatchTST), we
add three exogenous channels aligned with the daily target:
\begin{itemize}[leftmargin=*]
    \item \textbf{Temperature.} Daily mean temperature in Boulder
    pulled from the Open-Meteo historical weather
    API~\cite{openmeteo}. We picked mean over min/max because most
    charging happens during commute hours, when the mean is the
    closest proxy for the temperature actually seen by the chargers.
    \item \textbf{Day-of-week flags.} A binary \emph{is\_weekend}
    flag, plus the integer day of week as a feature when relevant. We
    do \emph{not} use a holiday-aware day-of-week encoding because
    holidays are handled separately.
    \item \textbf{Holiday indicator.} A binary flag for US public
    holidays, generated from the \texttt{holidays} Python package
    using the federal calendar plus Colorado state holidays.
\end{itemize}
These three variables are known at forecast time (day-of-week and
holiday from the calendar, temperature from a standard one-week
weather forecast), so feeding them as future exogenous inputs is
causally legitimate.

\subsection{Stationarity and Seasonality Analysis}

Before we picked any model we ran the Augmented Dickey-Fuller test on
the raw daily series. The ADF statistic is $-3.097$ with $p$-value
$0.027$ at 24 lags (AIC selection), which lets us reject the unit-root
null at the 5\% level. The series is stationary, not a random walk;
we therefore use $d=0$ as the default differencing order for the
(S)ARIMA family, with the search space allowing $d \in \{0,1\}$ for
SARIMA's first difference of the seasonally differenced series.

Figure~\ref{fig:acf} shows the autocorrelation function of the raw
series. The picture is unambiguous: clear peaks at lags 7, 14, 21, 28
and 35, with a slow decay in between, telling us the weekly cycle is
the dominant seasonal component. There is no annual peak visible in
this lag range because we only show 42 lags, but the seasonal
behaviour at the year scale is clear in the time-series plot
(Fig.~\ref{fig:series}, where summer months consistently sit above
the mean). The PACF (not shown for space reasons) has significant
spikes at lags 1 and 7, suggesting an AR(1) main component with a
seasonal AR at lag 7 --- which is the search shape we use for the
(S)ARIMA grid in Section~\ref{sec:models}.

\begin{figure}[!htbp]
\centering
\includegraphics[width=0.85\textwidth]{fig_acf}
\caption{Autocorrelation of the daily series. The repeating peaks at
multiples of 7 are the signature of a strong weekly seasonal cycle;
this is the structural reason classical (S)ARIMA models do not need
to be replaced by anything fancier to be competitive.}
\label{fig:acf}
\end{figure}
